\section{Methods}
\label{sec:methods}

\subsection{Policy families}
Table~\ref{tab:families} lists the four Level~1 families. All share the observation of
Section~\ref{sec:setup}, the normalized action, a masked mean-squared error in normalized
action space (the diffusion family's error is on the predicted noise), float32 training on
the CPU and one shared action decoder.
\begin{itemize}
\item \textbf{BC-MLP} maps the current observation to the next action ($T_o = T_p = T_a =
  1$), through hidden layers of width \rf{mlp.hidden}.
\item \textbf{BC-RNN} runs an LSTM~\citep{hochreiter1997lstm} with hidden size
  \rf{rnn.hidden} over the last \rf{hor.rnn} observations and predicts the next action from its last valid step. At every decision the
  LSTM restarts from a zero state over the window, exactly as in training, so no hidden state
  persists between decisions or episodes. This differs from robomimic's BC-RNN, which carries
  its hidden state for several steps; we chose exact train/inference consistency.
\item \textbf{BC-MLP-chunk} keeps BC-MLP's hidden layers but takes the diffusion policy's
  window of $T_o = \rf{hor.obs}$ observations and predicts $T_p = \rf{hor.pred}$ future
  actions, of which it executes the first $T_a = \rf{hor.exec}$ before observing again. It
  shares the diffusion policy's inputs, horizons and training budget and is the control for
  H2-c. It differs in the head (a deterministic MLP against a denoising U-Net evaluated with
  EMA weights), in the action bound described below, and in the tuning trial each family
  selected.
\item \textbf{Diffusion U-Net} is a conditional denoising diffusion model over the
  normalized $T_p \times \rf{dim.act}$ action chunk~\citep{chi2023diffusion}. The
  $T_o \times \rf{dim.obs}$ observation window is flattened into a global condition, with no
  learned observation encoder. A 1D temporal U-Net~\citep{janner2022diffuser} with down
  widths (\rf{diff.width1}, \rf{diff.width2}), GroupNorm~\citep{wu2018groupnorm} and
  Mish~\citep{misra2019mish} activations predicts the noise; every residual block receives a
  FiLM scale and bias~\citep{perez2018film} from the diffusion-timestep embedding and the
  condition. The architecture is adapted from robomimic's Diffusion Policy implementation,
  scaled down to match BC's parameter budget and trained with the same optimizer recipe as
  BC rather than robomimic's warm-up schedule; its scheduler semantics were checked against
  diffusers~\citep{diffusers}.
\end{itemize}

\begin{table}[t]
\centering\small
\caption{Level~1 policy families. Parameter counts are exact; all four lie within the budget
rule of the protocol (diffusion's count is not more than twice BC-MLP's, so no
capacity-matched BC variants were added).}
\label{tab:families}
\resizebox{\linewidth}{!}{% Generated by scripts/paper_assets.py (DECISIONS D-094). Do not edit.
% Sources (path sha256):
%   configs/models/diffusion_unet.yaml ca181c22a4af5499a67066fe2d78c06aeff5a889212ad3cceed7265fc8228bbf
%   reports/session4/param_counts.json 0ec7b26a0a329737cc0b18a968f5276d7130a5f66614305985505a16532486cb
\begin{tabular}{llll}
\toprule
Family & Network & Horizons & Parameters \\
\midrule
BC-MLP & MLP 2 $\times$ 1024, ReLU & $T_o = 1$, $T_p = 1$, $T_a = 1$ & 1,077,255 \\
BC-RNN & LSTM 2 $\times$ 296 & context 10, $T_p = 1$, $T_a = 1$ & 1,080,703 \\
BC-MLP-chunk & MLP 2 $\times$ 1024, ReLU & $T_o = 2$, $T_p = 16$, $T_a = 8$ & 1,204,336 \\
Diffusion U-Net & 1D U-Net (64, 128), FiLM; DDPM $K = 100$, DDIM-10, EMA & $T_o = 2$, $T_p = 16$, $T_a = 8$ & 1,092,551 \\
\bottomrule
\end{tabular}
}
\end{table}

\subsection{Diffusion training and sampling}
Training uses the DDPM forward process~\citep{ho2020ddpm} with $K = \rf{diff.T}$ steps and
the squared-cosine noise schedule~\citep{nichol2021improved}:
$A_t = \sqrt{\bar\alpha_t}\,A_0 + \sqrt{1-\bar\alpha_t}\,\epsilon$ for a normalized
demonstrated chunk $A_0$, a timestep $t$ drawn uniformly and Gaussian noise $\epsilon$. The
network predicts $\epsilon$ and the loss is the masked squared error over the valid
(unpadded) action steps. At inference the policy runs DDIM~\citep{song2021ddim} with
\rf{diff.ddim} steps and $\eta = 0$, so a chunk is a deterministic function of its initial
noise; the predicted clean chunk is clipped to $[-1, 1]$ at every step. In min--max
normalized space this keeps every diffusion action within the range of the training actions
on each dimension, a bound the BC heads do not have (their outputs are clipped only to the
controller limits; Section~\ref{sec:limitations}). The evaluated
network is an exponential moving average (EMA) of the weights with power \rf{diff.ema}; the
raw trained weights are kept in every checkpoint for ablation A1. DDIM with \rf{diff.ddim}
steps was fixed before any reported run because the robomimic default (DDPM with all
\rf{diff.T} steps) made every canonical evaluation several times longer; A2 compares the two.

\subsection{Optimization and matched budgets}
The base configuration of every family uses AdamW~\citep{loshchilov2019adamw} with learning
rate \rf{train.lr}, batch \rf{train.batch} and no gradient clipping; the tuning trials below
change only the optimizer. Every main run trains for
\rf{train.steps} gradient steps and saves a checkpoint every \rf{train.ckpt} steps
(\rf{train.K} checkpoints), identically for all families and both tasks, so every family
sees the same number of training samples. Training ran on the local CPU with one torch
thread per run.

\paragraph{Tuning budget.} Each family received exactly \rf{train.trials} configurations per
task, each trained with the first training seed at the main budget and declared before the first run: the
pilot configuration; a cosine learning-rate schedule annealed to zero; a constant learning
rate of \rf{train.lr3}; and weight decay \rf{train.wd4}. The architecture, horizons, batch
and budget never changed. The configuration whose selected checkpoint had the most
validation successes (ties: lower mean first-success step, then the lower trial number) was
then trained with the other two seeds. The chosen trials were, on Lift, \rf{sel.lift.bcmlp.chosen}
(BC-MLP), \rf{sel.lift.bcrnn.chosen} (BC-RNN), \rf{sel.lift.chunk.chosen} (BC-MLP-chunk) and
\rf{sel.lift.diff.chosen} (diffusion), and on Can \rf{sel.can.bcmlp.chosen},
\rf{sel.can.bcrnn.chosen}, \rf{sel.can.chunk.chosen} and \rf{sel.can.diff.chosen}.

\subsection{Observation-modality extensions}
Two further studies replace the privileged state by sensor observations on Lift, each with
its own freeze. Both keep the diffusion head, horizons and action encoding above, train for
\rf{vis.steps} steps with batch \rf{vis.batch} and \rf{vis.K} checkpoints on one NVIDIA L4
GPU with deterministic CUDA kernels, and evaluate on the same CPU platform as Level~1.
\begin{itemize}
\item \textbf{RGB study.} Two cameras (the front \code{agentview} and the wrist
  \code{robot0_eye_in_hand}) at \rf{vis.image}$\times$\rf{vis.image} pixels, rendered from
  the recorded demonstration states, each randomly cropped during training (centre crop at evaluation) and encoded by a
  ResNet-18~\citep{he2016resnet} with GroupNorm and random initialization followed by a
  spatial soft-argmax~\citep{levine2016visuomotor},
  plus \rf{pc.proprio}-dimensional proprioception (end-effector pose and gripper; no object
  state). Visual diffusion is compared with visual BC-chunk, which has the same encoder and
  horizons and a deterministic head.
\item \textbf{Level~2 point clouds (DP3-inspired).} One world-frame cloud per step, fused
  from the depth images of the same two cameras, cropped to a fixed workspace box and reduced
  to \rf{pc.points} points by deterministic farthest-point sampling. A per-point MLP
  (widths \rf{pc.enc.w1} to \rf{pc.enc.w3}, LayerNorm, ReLU) with max pooling and a
  \rf{pc.enc.out}-dimensional projection encodes it, following the idea of
  DP3~\citep{ze2024dp3} and PointNet~\citep{qi2017pointnet}; the same proprioception is
  added. The cloud uses only depth, camera intrinsics and camera poses (the wrist camera's
  from robot kinematics). No segmentation, object pose or simulator label enters it, so the
  arm is sensor-like rather than privileged. The depth is noise-free and the calibration is
  exact. The reference arm is the RGB study's diffusion arm, frozen and never retrained.
\end{itemize}
